\documentclass[10pt,a4paper]{amsart}
\usepackage[margin=19mm]{geometry}
\usepackage{amsmath,amssymb,mathtools}
\usepackage[scr=rsfs,cal=euler]{mathalfa}
\usepackage{xfrac}
\usepackage[unicode,hidelinks,bookmarks=false]{hyperref}
\newcommand{\ie}{i.e.\ }
\newcommand{\cf}{cf.\ }
\newcommand{\resp}{respectively\ }
\newcommand{\Subsection}{\subsection}
\providecommand{\nobreakdash}{\nobreak\hbox{-}}
\setlength{\emergencystretch}{2em}
\multlinegap0pt
\usepackage{bm}
\usepackage{bbm}
\usepackage{dsfont}
\usepackage{xspace}
\usepackage{cancel}
\usepackage{mathtools}
\usepackage{cleveref}
\usepackage[shortlabels]{enumitem}
\usepackage{amsthm}
\makeatletter
\@ifundefined{definition}{\newtheorem{definition}{Definition}}{}
\@ifundefined{lemma}{\newtheorem{lemma}[definition]{Lemma}}{}
\@ifundefined{theorem}{\newtheorem{theorem}[definition]{Theorem}}{}
\makeatother
\DeclareMathOperator{\maxmat}{N_{\mathrm{max}}}
\DeclareMathOperator{\minmat}{N_{\mathrm{min}}}
\newcommand{\maxedges}{E_{\mathrm{max}}}
\newcommand{\minedges}{E_{\mathrm{min}}}
\renewcommand{\subset}{\subseteq}
\title[Temporal Cliques Admit Linear Spanners]{Temporal Cliques Admit Linear Spanners}
\author{Julia Baligacs}
\date{Source: arXiv:2606.05156v1 (CC BY 4.0)}
\begin{document}
\maketitle

\noindent\emph{Source and licence.} Temporal Cliques Admit Linear Spanners. Version arXiv:2606.05156v1 (CC BY 4.0), distributed under the Creative Commons Attribution 4.0 International licence, as recorded in the arXiv licence metadata for this exact version and shown on the abstract page https://arxiv.org/abs/2606.05156v1. Full rights notice: licenses/arxiv-2606.05156-CC-BY-4.0.txt.\par\smallskip

\noindent\textit{Editorial scope.} Counted unit: the Theorem whose source label is \texttt{thm:mainthm\_bicliques\_stronger} in the article (source0.tex lines 894--898, source label \texttt{thm:mainthm\_bicliques\_stronger}), delivered together with the article's own complete original proof of it (source0.tex lines 901--936, one \texttt{proof} environment of 36 lines). No mathematics is written, repaired or re-derived: statement and proof are the article's own text line for line. Required context retained uncounted: lem:dismountable (lines 451--454); lem:final (lines 865--874). Every definition, display and label of those lines is retained unchanged. Omitted from this package and not redistributed: 1--450, 455--864, 875--893, 899--900, 937--1238. Nothing inside a retained excerpt is omitted or altered: every retained body is delivered exactly as the article's lines are written, so no omission receipt of any kind accompanies this package. Citations: the retained excerpts carry no citation command at all, so no bibliography entry of the article is transcribed and no reference is redistributed. Reference adaptation: none was needed and none was made; every reference inside the delivered unit resolves inside it, so no unresolved reference or citation is produced. Printed slips kept literal: no printed word, symbol, hypothesis or number of the counted statement or of its proof was altered. Layout and class: the article's own class is \texttt{article} with packages including lineno, fontenc, lmodern, amsmath, amssymb, amsthm, dsfont, mathtools, thmtools, thm-restate, xspace, graphicx, tikz, xcolor; the wrapper is the lane's pinned \texttt{amsart} editorial wrapper at 10pt on A4, which restates the theorem-like environments the retained text needs and adds the article's own macro definitions that the retained text needs (\texttt{\textbackslash maxedges}, \texttt{\textbackslash maxmat}, \texttt{\textbackslash minedges}, \texttt{\textbackslash minmat}, \texttt{\textbackslash subset}), with extra wrapper packages (bm, bbm, dsfont, xspace, cancel, mathtools, cleveref, enumitem). The delivered numbering is the unit's own. Assets: no retained span contains a figure environment, a graphics-inclusion command, a PDF-page inclusion, a table or a TikZ picture; no image file is copied into this package, the package declares no asset dependency and no image file is redistributed with it. Licence: version arXiv:2606.05156v1 of this article is distributed under the Creative Commons Attribution 4.0 International licence, as recorded in the arXiv licence metadata for this exact version and shown on the abstract page https://arxiv.org/abs/2606.05156v1. A full rights notice is delivered at licenses/arxiv-2606.05156-CC-BY-4.0.txt. Characters: every retained excerpt is pure ASCII, so the delivered engine, XeLaTeX with its default Latin Modern text font, reports no missing character.
\begin{lemma}[Dismountability]
\label{lem:dismountable}
Given a temporal bi-clique $G=(S,T,\lambda)$, there exist $S'\subseteq S$, $T'\subseteq T$, and a set $E'$ of at most $2(|S|+|T|-|S'|-|T'|)$ edges such that $G[S',T']$ is extremally matched and, if $E^*$ is a temporal spanner of $G[S',T']$, then $E^*\cup E'$ is a temporal spanner of $G$.
\end{lemma}

\begin{lemma}
\label{lem:final}
Fix $s^*\in S$, consider the $s^*$-ordered labeling, and set $k:=\lfloor n/2 \rfloor$.
Let $S':=\{s_0, \dots, s_k\}$ and $T':=\{t_k, \dots, t_{n-1}\}$.
Then there exists a set $E^*$ of at most $6n$ edges containing $\minedges \cup \maxedges$ that covers either
\begin{enumerate}[label=(\roman*)]
  \item \label{lem:final_cond_source} $(S,T')$, or
  \item \label{lem:final_cond_target} $(S',T)$.
\end{enumerate}
\end{lemma}

\begin{theorem}
\label{thm:mainthm_bicliques_stronger}
Every extremally matched temporal bi-clique $(S,T,\lambda)$ with $n:=|S|=|T|$ contains a temporal spanner of size $14n$.
Moreover, this spanner also covers $(S\cup T, S\cup T)$.
\end{theorem}

\begin{proof}
Let $f(n)$ denote the maximum size of a lightest spanner of an extremally matched temporal bi-clique of size $n$, that is,
\begin{equation*}
f(n):=\max \{\min\{|E'|: E' \text{ is a spanner of } G\}: G=(S,T,\lambda) \text{ is extremally matched with } |S|=n\}.
\end{equation*}
Let $G_n=(S,T,\lambda)$ be an extremally matched temporal bi-clique realizing this maximum.
Fix~${s^*\in S}$, consider the $s^*$-ordered labeling, and let $S',T'$ be as in \cref{lem:final}.
By \cref{lem:final}, there exists a set $E^*$ of at most $6n$ edges covering either $(S,T')$ or $(S',T)$.
In the first case, it remains to find a spanner of $(S,T\setminus T')$, and in the second a spanner of $(S\setminus S',T)$.
In both cases, call this remaining subgraph $G'$.

Note that $|S \setminus S'| = |\{s_{\lfloor n/2 \rfloor +1}, \dots, s_{n-1}\}| \leq \lfloor n/2 \rfloor$, and similarly $|T \setminus T'| = |\{t_0, \ldots, t_{\lfloor n/2 \rfloor -1}\}| = \lfloor n/2 \rfloor$.
Hence, in both cases, $G'$ is a bi-clique with at most $\lfloor 3n/2 \rfloor$ vertices in total, with the smaller side having size at most $\lfloor n/2 \rfloor$ and the larger side having size $n$.
Applying dismountability, i.e., \cref{lem:dismountable} to $G'$, we obtain $S''\subset S$ and $T''\subset T$ such that $G[S'',T'']$ is extremally matched with~$|S''|=|T''|\leq \lfloor n/2 \rfloor$, and a set of edges $E''$ such that adding $E''$ to any spanner of $G[S'',T'']$ yields a spanner of $G'$, where
\begin{equation*}
|E''|\leq 2\cdot (\lfloor 3n/2 \rfloor - |S''|-|T''|)
\overset{|S''|=|T''|}{\leq}3n - 4|S''|.
\end{equation*}
To summarize, we have selected the edges $E^*\cup E''$ and it only remains to find a spanner of $G[S'',T'']$ to obtain a spanner of $G_n$.
Setting $x_n := |S''| \leq \lfloor n/2 \rfloor$, we have
\begin{equation*}
f(n)\leq |E^*|+|E''|+f(|S''|)
\leq
6n + 3n - 4x_n + f(x_n)
=9n- 4x_n + f(x_n).
\end{equation*}
We now prove by strong induction that $f(n)\leq 14n$. The base case $f(1)=1$ is immediate (and we set $f(0):=0$ for the case where $x_n=0$, i.e., the recursion ends). For the induction step, we have
\begin{equation*}
f(n)\leq 9n-4x_n+f(x_n)
\overset{\text{ind.~hyp.}}{\leq} 9n + (14-4)x_n 
\overset{x_n\leq n/2}{\leq}
 9n + 10 \cdot \frac{n}{2} \leq 14n.
\end{equation*}
Hence spanners of the claimed size exist.
Last, the final edge set contains $\minedges \cup \maxedges$ as it is contained in the edge set added in the first recursion level chosen according to \cref{lem:final}. This ensures that the spanner does in fact preserve all temporal reachabilities: any target~$t$ reaches the same set of vertices as $\minmat(t)$, since prefixing any temporal walk beginning in $\minmat(t)$ with the edge $\{t, \minmat(t)\}$ yields a valid temporal walk (where we have used that $G$ is extremally matched). Symmetrically, any source~$s$ is reachable from the same set of vertices that reach $\maxmat(s)$, since appending $\{s, \maxmat(s)\}$ to any temporal walk ending at $\maxmat(s)$ again yields a temporal walk.
\end{proof}

\end{document}
